%--- Referencias: cabecera + nota de fechas (reescrita) ---
\newpage
\section*{Referencias bibliográficas}
\addcontentsline{toc}{section}{Referencias bibliográficas}
\markboth{Referencias bibliográficas}{Referencias bibliográficas}

Todas las fuentes que aparecen en este trabajo son \textbf{documentación oficial del fabricante o de la plataforma} (Android Developers, Google Play Console, Apple, Material Design, W3C, Bluetooth SIG).
\begin{warnbox}
\textbf{\color{accent!80!black}\faExclamationTriangle\ \ Las fechas de esta parte hay que revisarlas.}
Los datos de versiones que aparecen en el trabajo (niveles de API, políticas de Google Play,
requisitos de Android Studio, umbrales de memoria) cambian con frecuencia. Las páginas que cito
están consultadas a fecha de \textbf{\texttt{FECHA DE CONSULTA}}, y antes de entregar las vuelvo
a abrir para que las cifras del apartado 1 y del apartado 4 sigan siendo correctas. Las cuatro que
más se mueven son estas:
\begin{itemize}[leftmargin=*,itemsep=0.15em,topsep=0.2em]
  \item \href{https://developer.android.com/tools/releases/platforms}{\emph{Android API levels}}
  \item \href{https://support.google.com/googleplay/android-developer/answer/11926878}{\emph{Target API level requirements} (Google Play)}
  \item \href{https://developer.android.com/studio}{\emph{Android Studio} (versión y requisitos)}
  \item \href{https://developer.android.com/topic/performance/vitals}{\emph{Android vitals} (umbrales de memoria)}
\end{itemize}
\end{warnbox}
%--- Lista de referencias + mapa de fuentes (verbatim, sin el aviso final) ---
\subsection*{Fuentes de Android Developers (documentación técnica)}
{\small
\begin{itemize}[leftmargin=*,itemsep=0.1em,label={}]
  \item Android Developers. (2026). \emph{Android vitals}. \url{https://developer.android.com/topic/performance/vitals}
  \item Android Developers. (2026). \emph{Support different screen densities}. \url{https://developer.android.com/training/multiscreen/screendensities}
  \item Android Developers. (2026). \emph{Screen sizes}. \url{https://developer.android.com/training/multiscreen/screensizes}
  \item Android Developers. (2026). \emph{Configuration.screenLayout}. \url{https://developer.android.com/reference/android/content/res/Configuration}
  \item Android Developers. (2026). \emph{Android API levels}. \url{https://developer.android.com/tools/releases/platforms}
  \item Android Developers. (2026). \emph{uses-sdk element}. \url{https://developer.android.com/guide/topics/manifest/uses-sdk-element}
  \item Android Developers. (2026). \emph{Android 17 is here}. \url{https://developer.android.com/blog/posts/android-17-is-here}
  \item Android Developers. (2026). \emph{Behavior changes: all apps}. \url{https://developer.android.com/about/versions/17/behavior-changes-all}
  \item Android Developers. (2026). \emph{Android 17 release notes}. \url{https://developer.android.com/about/versions/17/release-notes}
  \item Android Developers. (2026). \emph{Storage overview}. \url{https://developer.android.com/training/data-storage}
  \item Android Developers. (2026). \emph{APK Analyzer}. \url{https://developer.android.com/studio/debug/apk-analyzer}
  \item Android Developers. (2026). \emph{Android App Bundles}. \url{https://developer.android.com/guide/app-bundle/dynamic-delivery}
  \item Android Developers. (2026). \emph{ActivityManager}. \url{https://developer.android.com/reference/android/app/ActivityManager}
  \item Android Developers. (2026). \emph{Guide to Android app architecture}. \url{https://developer.android.com/topic/libraries/architecture}
  \item Android Developers. (2026). \emph{5G}. \url{https://developer.android.com/develop/connectivity/5g}
  \item Android Developers. (2026). \emph{Bluetooth}. \url{https://developer.android.com/develop/connectivity/bluetooth}
  \item Android Developers. (2026). \emph{Bluetooth LE Audio overview}. \url{https://developer.android.com/develop/connectivity/bluetooth/ble-audio/overview}
  \item Android Developers. (2026). \emph{Near Field Communication}. \url{https://developer.android.com/develop/connectivity/nfc}
  \item Android Developers. (2026). \emph{Wi-Fi}. \url{https://developer.android.com/develop/connectivity/wifi}
  \item Android Studio Community. (2026). \emph{Wi-Fi 7}. \url{https://source.android.com/docs/core/connect/wifi-7}
\end{itemize}
}

\subsection*{Fuentes de Android Studio (Apartado 3)}
{\small
\begin{itemize}[leftmargin=*,itemsep=0.1em,label={}]
  \item Android Developers. (2026). \emph{Android Studio}. \url{https://developer.android.com/studio}
  \item Android Developers. (2026). \emph{Install Android Studio}. \url{https://developer.android.com/studio/install}
  \item Android Developers. (2026). \emph{Create a new Android Studio project}. \url{https://developer.android.com/studio/projects/create-project}
  \item Android Developers. (2026). \emph{Android Studio project templates}. \url{https://developer.android.com/studio/projects/templates}
  \item Android Developers. (2026). \emph{Run the Android Emulator}. \url{https://developer.android.com/studio/run/emulator}
  \item Android Developers. (2026). \emph{Configure hardware acceleration}. \url{https://developer.android.com/studio/run/emulator-acceleration}
  \item Android Developers. (2026). \emph{Create and manage Android virtual devices}. \url{https://developer.android.com/studio/run/managing-avds}
  \item Android Developers. (2026). \emph{Android Emulator troubleshooting}. \url{https://developer.android.com/studio/run/emulator-troubleshooting}
  \item Android Developers. (2026). \emph{Android Studio troubleshooting}. \url{https://developer.android.com/studio/troubleshoot}
  \item Android Developers. (2026). \emph{Android UI development is Compose-first} [blog, 19/05/2026]. \url{https://developer.android.com/blog/posts/android-ui-development-is-compose-first}
\end{itemize}
}

\subsection*{Fuentes de accesibilidad y diseño}
{\small
\begin{itemize}[leftmargin=*,itemsep=0.1em,label={}]
  \item Android Developers. (2026). \emph{Make your app more accessible}. \url{https://developer.android.com/guide/topics/ui/accessibility/apps}
  \item Android Developers. (2026). \emph{Principles for improving app accessibility}. \url{https://developer.android.com/guide/topics/ui/accessibility/principles}
  \item Android Developers. (2026). \emph{Accessibility in Jetpack Compose}. \url{https://developer.android.com/develop/ui/compose/accessibility}
  \item Android Developers. (2026). \emph{TalkBack y accesibilidad en Android}. \url{https://support.google.com/accessibility/android/answer/6006564}
  \item Android Developers. (2026). \emph{Use TalkBack}. \url{https://support.google.com/accessibility/android/answer/6007100}
  \item Android Developers. (2026). \emph{Use braille with TalkBack}. \url{https://support.google.com/accessibility/android/answer/16323943}
  \item Android Developers. (2026). \emph{Change font size and display size}. \url{https://support.google.com/accessibility/android/answer/11183305}
  \item Google. (2026). \emph{Target API level requirements}. \url{https://support.google.com/googleplay/android-developer/answer/11926878}
\end{itemize}
}

\subsection*{Fuentes de estándares de diseño y accesibilidad}
{\small
\begin{itemize}[leftmargin=*,itemsep=0.1em,label={}]
  \item World Wide Web Consortium (W3C). (2023). \emph{Web Content Accessibility Guidelines (WCAG) 2.2}. \url{https://www.w3.org/TR/WCAG22/}
  \item Material Design 3. (2026). \emph{Color contrast}. \url{https://m3.material.io/foundations/designing/color-contrast}
  \item Material Design 3. (2026). \emph{Structure}. \url{https://m3.material.io/foundations/designing/structure}
  \item Material Design 3. (2026). \emph{Grids and spacing}. \url{https://m3.material.io/foundations/layout/grids-spacing/density}
\end{itemize}
}

\subsection*{Fuentes de fabricantes}
{\small
\begin{itemize}[leftmargin=*,itemsep=0.1em,label={}]
  \item Samsung Electronics. (2026). \emph{Galaxy S26 Ultra --- Especificaciones técnicas}. \url{https://www.samsung.com/}
  \item Google. (2026). \emph{Pixel 10 Pro XL --- Especificaciones técnicas}. \url{https://store.google.com/}
  \item OnePlus. (2026). \emph{OnePlus 13 --- Especificaciones}. \url{https://www.oneplus.com/}
  \item Apple Inc. (2026). \emph{iPhone 17 Pro Max --- Especificaciones técnicas}. \url{https://www.apple.com/iphone-17-pro-max/technical-specs/}
  \item Motorola. (2026). \emph{moto g47 --- Especificaciones técnicas}. \url{https://www.motorola.com/}
  \item Apple Inc. (2026). \emph{Apple Security Releases}. \url{https://support.apple.com/en-us/148282}
\end{itemize}
}

\subsection*{Fuentes de lenguajes y frameworks del Apartado 2}
{\small
\begin{itemize}[leftmargin=*,itemsep=0.1em,label={}]
  \item Kotlin Foundation. (2026). \emph{Kotlin documentation}. \url{https://kotlinlang.org/docs/home.html}
  \item Oracle. (2026). \emph{Java Documentation}. \url{https://docs.oracle.com/en/java/}
  \item Flutter. (2026). \emph{Flutter documentation}. \url{https://docs.flutter.dev/}
  \item Meta Open Source. (2026). \emph{React Native documentation}. \url{https://reactnative.dev/docs/getting-started}
  \item Microsoft. (2026). \emph{Build mobile apps with .NET MAUI}. \url{https://learn.microsoft.com/dotnet/maui/}
  \item Apple Inc. (2026). \emph{Swift Documentation}. \url{https://developer.apple.com/documentation/swift/}
\end{itemize}
}

\subsection*{Mapa de fuentes por apartado}

{\small
\begin{table}[htbp]
\centering
\begin{tabularx}{\textwidth}{@{}>{\RaggedRight\arraybackslash}p{3.5cm} >{\RaggedRight\arraybackslash}X@{}}
\toprule
\rowcolor{primary}
\textcolor{white}{\sffamily\bfseries Apartado} & \textcolor{white}{\sffamily\bfseries Fuentes principales} \\
\midrule
\textbf{1. Limitaciones} & Android vitals · \emph{Guide to Android app architecture} · Bluetooth LE Audio · fichas técnicas de los cinco fabricantes \\
\rowcolor{soft}
\textbf{2. Tecnologías} & Documentación de Kotlin, Oracle, Flutter, React Native, Microsoft y Apple · blog Compose (19/05/2026) \\
\textbf{3. Entorno} & Android Studio · Install · Create project · Emulator · Hardware acceleration · AVDs · Troubleshooting · uses-sdk \\
\rowcolor{soft}
\textbf{4. Config. dispositivos} & Screen densities · Screen sizes · Configuration.screenLayout · API levels · Target API · Storage · App Bundles · APK Analyzer · vitals · ActivityManager · 5G · Bluetooth · NFC · Wi-Fi · Wi-Fi 7 · Behavior changes \\
\textbf{5. Perfiles} & Make your app more accessible · Principles · Accessibility in Compose · TalkBack · WCAG 2.2 · Material Design 3 · Bluetooth LE Audio \\
\bottomrule
\end{tabularx}
\caption{Mapa de fuentes por apartado}
\end{table}
}
